\paragraph{Data and task.} A dataset contains $S$ subjects. Subject $s$ has paired sequences $\{(x_{s,t}, y_{s,t})\}_{t=1}^{T_s}$, where $y_{s,t}\in\R^{V}$ is the fMRI volume acquired at time $t$ ($V = 64\times 64\times 30$ voxels) and $x_{s,t}\in\R^{C\times F\times W}$ is the short-time Fourier magnitude of the $C$ EEG channels over the $W$ seconds preceding it (so that it covers the haemodynamic delay). A synthesizer is a conditional predictive distribution $p_\theta(y\mid x)$ trained on a set of training subjects $\mathcal{S}_{\mathrm{tr}}$ and evaluated on disjoint test subjects $\mathcal{S}_{\mathrm{te}}$. A disjoint set of calibration subjects $\mathcal{S}_{\mathrm{cal}}$ is used for conformal calibration. All splits are by subject; no subject contributes volumes to two sets.

\paragraph{Scores.} For a test pair $(x,y)$ and a predictive distribution $p$ we use the voxel-averaged negative log-likelihood $\mathrm{NLL} = -\frac{1}{|B|}\sum_{v\in B}\log p(y_v\mid x)$ and the continuous ranked probability score $\mathrm{CRPS} = \frac{1}{|B|}\sum_{v\in B}\E|Y_v-y_v| - \tfrac12\E|Y_v-Y'_v|$ with $Y,Y'\sim p(\cdot\mid x)$ independent \citep{gneiting2007strictly}, both over a fixed brain mask $B$. We also use the brain RMSE of the predictive mean, the prediction interval coverage probability (PICP) and width of the central $\alpha$-intervals, Spearman's $\rho$ between the predictive standard deviation and the absolute error, and the area under the sparsification error curve (AUSE). The mask $B$ is the Otsu threshold \citep{otsu1979threshold} of the training subjects' mean volume, so it depends on no test data. Every score is first averaged over each test subject's volumes, and then over seeds; subjects, not volumes, are the unit of every test.

\paragraph{The EEG-conditional question.} The checks above evaluate $p_\theta(y\mid x)$ against $y$. They do not ask whether $p_\theta$ depends on $x$. We say that a synthesizer's predictive distribution is \emph{informative about the EEG} if it passes two comparisons.

\emph{(i) EEG-free references.} It must score better than predictive distributions that never see EEG:
\begin{itemize}
	\item the \textbf{population template} $p_0(y) = \prod_v \mathcal{N}(y_v; \bar m_v, \bar s_v^2)$, with the voxelwise mean and standard deviation of all training volumes;
	\item the \textbf{stimulus-locked template} $p_{\mathrm{lock}}(y\mid t) = \prod_v\mathcal{N}(y_v; \bar m_{v,t}, \bar s^2_{v,t})$, with moments over training subjects at the same volume index $t$. In a block design with the same timing for every subject, this template predicts the task-evoked response without any EEG;
	\item the \textbf{own-subject oracle} $p_{\mathrm{own}}(y) = \prod_v \mathcal{N}(y_v; m^{(s)}_v, (s^{(s)}_v)^2)$, with the moments of the test subject's own volumes. It uses test data and is an upper bound for any predictor that gets the subject's static image right and nothing else.
\end{itemize}
A method passes if its subject-level NLL and CRPS are lower than those of $p_0$ and $p_{\mathrm{lock}}$, by a one-sided Wilcoxon signed-rank test over test subjects \citep{wilcoxon1945individual}.

\emph{(ii) EEG nulls.} The same trained model is evaluated again with its input replaced by a null $\tilde x$ that keeps the marginal statistics of the EEG but breaks its pairing with the fMRI:
\begin{itemize}
	\item a \textbf{shift null}, the subject's EEG circularly shifted by at least 30 volumes, which preserves its spectrum and slow trends;
	\item a \textbf{swap null}, the EEG of another test subject at the same volume index, which preserves anything locked to the shared task timing and the run (drifts, block design), but not the subject.
\end{itemize}
A method uses the EEG if a score of the real pairing is better than that of each null. The scores are the median voxel temporal correlation $r$ between predicted and observed time courses, the pooled detrended correlation, the centred $R^2$ of the voxel time courses, and the NLL. Beating the shift null but not the swap null means the model has learned something run- or task-locked, not subject-specific EEG information.

\paragraph{Coverage at the subject level.} Split conformal intervals $[\mu_v - q\,\sigma_v, \mu_v + q\,\sigma_v]$ are calibrated on $\mathcal{S}_{\mathrm{cal}}$ with the normalized score $|y_v-\mu_v|/\sigma_v$. Here $\mu,\sigma$ are the predictive mean and standard deviation, and $\sigma\equiv 1$ for a point predictor. We calibrate either one quantile for all voxels (global) or one per voxel. The quantile regression head uses the CQR score instead \citep{romano2019conformalized}. Conformal coverage holds for exchangeable units \citep{lei2018distribution}. Here the units are subjects, not volumes or voxels, and with two calibration subjects per fold the guarantee does not hold for a single test subject. We therefore report the empirical coverage on held-out subjects instead of the nominal value.
